% -*- mode: latex; coding: utf-8 -*-
%-----------------------------------------------------------------------------
% pxTAB -- User manual and API reference
% Part of P3CTeX. Author: Scvria.
%
% Build from directory tex/ (TEXINPUTS so package is found):
%   set TEXINPUTS=.;./latex;./code;
%   pdflatex -interaction=nonstopmode doc/pxTAB.tex
%   makeindex -s gind.ist -o pxTAB.ind pxTAB.idx
%   pdflatex -interaction=nonstopmode doc/pxTAB.tex
%-----------------------------------------------------------------------------
\documentclass[11pt,a4paper]{article}
\PassOptionsToPackage{colorlinks,urlcolor=blue,linkcolor=blue,hyperindex}{hyperref}
\usepackage[T1]{fontenc}
\usepackage[utf8]{inputenc}
\usepackage{lmodern}
\usepackage{amsmath}
\usepackage{pxTAB}

\title{The \texttt{pxTAB} package\\
  \large Predefined table styles and list-driven tables for \LaTeX2e}
\author{Scvria\\P3CTeX / UOC PEC Repository}
\date{2026/03/02 -- Version 0.1}

\begin{document}
\maketitle

\begin{abstract}
  \texttt{pxTAB} provides a small \LaTeX2e\ API to build tables from
  comma-separated rows or list-style data without writing raw
  \texttt{tabular} environments.  It offers a handful of predefined
  visual styles (plain, header, striped, boxed, grid) and a set of
  composed header styles (header-boxed, header-grid, header-striped),
  together with key--value options for consistent layout control.
\end{abstract}

\tableofcontents

\section{Loading the package}

\begin{verbatim}
\usepackage{pxTAB}
\end{verbatim}

Package options are forwarded to the internal key family \texttt{pxTAB}
and are equivalent to calling \verb|\pxTABsetup| in the preamble:

\begin{verbatim}
\usepackage[style=header,headerbgcolor=blue!15]{pxTAB}
\end{verbatim}

\section{Configuration (keys)}

The \texttt{pxTAB} key family controls both visual style and parsing
behaviour.  Keys can be set at load time or via:

\begin{verbatim}
\pxTABsetup{<key>=<value>,...}
\end{verbatim}

Defaults are chosen to be reasonable for course work and reports.
All keys affect subsequent tables until changed again (usual \LaTeX\
scoping rules apply).

\subsection*{Style and colours}

\begin{itemize}
  \item \texttt{style} (choice, default \texttt{plain})\\
    Selects the high-level style:
    \texttt{plain}, \texttt{header}, \texttt{striped},
    \texttt{boxed}, \texttt{grid}, or one of the composed
    header styles \texttt{header-boxed}, \texttt{header-grid},
    \texttt{header-striped}.  The composed styles behave like
    the corresponding base style (boxed, grid, striped) with
    an additional coloured header row.

  \item \texttt{headerbgcolor} (token list, default \texttt{gray!25})\\
    Background colour of the header row in \texttt{header} style.

  \item \texttt{headertextcolor} (token list, default \texttt{black})\\
    Text colour used in the header row.

  \item \texttt{altrowbgcolor} (token list, default \texttt{gray!10})\\
    Background colour for alternating rows in \texttt{striped} style.

  \item \texttt{altrowcolor} (meta)\\
    Alias for \texttt{altrowbgcolor}.

  \item \texttt{altrowtextcolor} (token list, default empty)\\
    Text colour used on alternating rows in \texttt{striped} style.

  \item \texttt{headerfont} (token list, default \verb|\bfseries|)\\
    Font-switch applied to the header row (e.g.\ \verb|\sffamily\bfseries|).

  \item \texttt{bodyfont} (token list, default empty)\\
    Font-switch applied to body rows.

  \item \texttt{rulecolor} (token list, default \texttt{black!55})\\
    Colour used for horizontal and vertical rules
    (applied via \verb|\arrayrulecolor|).

  \item \texttt{rulewidth} (dimension, default \texttt{0.4pt})\\
    Thickness of rules (applied via \verb|\arrayrulewidth|).
\end{itemize}

\subsection*{Geometry and layout}

\begin{itemize}
  \item \texttt{tabcolsep} (dimension, default \texttt{6pt})\\
    Column separation (mirrors the standard \verb|\tabcolsep| parameter).

  \item \texttt{arraystretch} (token list, default \texttt{1.2})\\
    Row stretch applied via \verb|\arraystretch|.

  \item \texttt{rowstretch} (meta)\\
    Alias for \texttt{arraystretch}.

  \item \texttt{width} (code, default \texttt{natural})\\
    Width of the table.  The default \texttt{natural} keeps the table
    at its natural (fit-content) width; tables are \emph{not} stretched
    to page width unless you explicitly set \texttt{width} or
    \texttt{tabwidth}.  The special values \texttt{linewidth} and
    \texttt{line-width} are treated as \verb|\linewidth|; any other
    value is used as-is (e.g.\ \verb|.8\linewidth| or \verb|120pt|).

  \item \texttt{tabwidth} (meta)\\
    Alias for \texttt{width}.

  \item \texttt{centering} (boolean, default \texttt{false})\\
    When \texttt{true}, inserts a \verb|\centering| before the internal
    \texttt{tabular} or \texttt{tabular*}.  This only changes the
    horizontal alignment of the table block (in a paragraph or inside a
    \texttt{table} float); it does not modify the computed width.  In
    particular, \verb|width=\linewidth| still produces a full-width
    table, but the block is explicitly centred.  The key can be set
    globally with \verb|\pxTABsetup| or per table in the optional
    \verb|[<options>]| argument of the commands below.

  \item \texttt{fontsize} (token list, default empty)\\
    Optional font-size command (such as \verb|\small| or \verb|\footnotesize|)
    applied to the table.

  \item \texttt{size} (meta)\\
    Alias for \texttt{fontsize}.

  \item \texttt{align} (choice, default \texttt{c})\\
    When \texttt{colspec} is not provided, controls automatic column
    alignment: \texttt{l}, \texttt{c}, or \texttt{r}.

  \item \texttt{alignment} (meta)\\
    Alias for \texttt{align}.

  \item \texttt{colspec} (token list, default empty)\\
    Explicit \texttt{tabular} column specification (e.g.\ \texttt{ll},
    \texttt{|c|c|}).  When omitted, \texttt{pxTAB} builds one
    automatically from \texttt{align}.  Simple \texttt{l}/\texttt{c}/\texttt{r}
    specifications are checked against the inferred column count; on
    mismatch a warning is issued and an automatic specification is used.

\end{itemize}

\subsection*{Floats and captions}

The following keys enable a lightweight float wrapper around any
\texttt{pxTAB} table command.  They can be set at package level or
passed in the per-table options:

\begin{itemize}
  \item \texttt{caption} (token list, default empty)\\
    When non-empty, the table body is wrapped in a \texttt{table}
    environment and \verb|\caption{<caption>}| is emitted above the
    tabular.

  \item \texttt{label} (token list, default empty)\\
    Optional label placed immediately after the caption (or, if there is
    no caption, at the top of the float body).  A recommended convention
    is to use prefixes such as \texttt{tab:grades} so that
    \verb|\ref{tab:grades}| yields the table number in the usual way.

  \item \texttt{float} (token list, default empty)\\
    Placement specifier forwarded to the \texttt{table} environment
    (e.g.\ \texttt{htbp}).  When empty, \LaTeX\ uses its default float
    placement rules.
\end{itemize}

Any non-empty combination of \texttt{caption}, \texttt{label}, or
\texttt{float} activates the float wrapper; otherwise tables are
generated in-place without a surrounding \texttt{table} environment.
When both \texttt{centering=true} and float keys are used, the
\verb|\centering| is inserted inside the \texttt{table} environment,
before the caption and the tabular body.

\subsection*{Parsing}

\begin{itemize}
  \item \texttt{rowsep} (token list, default \texttt{;})\\
    Separator between rows when using list-style data.

  \item \texttt{cellsep} (token list, default \texttt{,})\\
    Separator between cells inside each row.

  \item \texttt{probe} (boolean, default \texttt{false})\\
    When \texttt{true}, writes diagnostic information to the \LaTeX\ log
    file during table rendering (style name, colspec, visual settings,
    etc.).  Intended for debugging and development; no effect on
    output.  QA harnesses may override the internal probe hook for
    automated assertions.

  \item \texttt{unknown} (code)\\
    Unknown keys are currently ignored silently (reserved for future
    extension).
\end{itemize}

\section{Styles by example}

This section illustrates the built-in styles using a small example
dataset.  In all cases we use:

\begin{verbatim}
\pxTABsetup{tabwidth=.8\linewidth}
\end{verbatim}

\subsection{Plain}

\begin{verbatim}
\pxTABsetup{style=plain,tabwidth=.8\linewidth}
\pxTable{Name, Age, City}
        {Alice, 20, Barcelona}
        {Bob, 22, Madrid}
\end{verbatim}

\pxTABsetup{style=plain,tabwidth=.8\linewidth}
\pxTable{Name, Age, City}
        {Alice, 20, Barcelona}
        {Bob, 22, Madrid}

\subsection{Header}

\begin{verbatim}
\pxTABsetup
  {
    style=header,
    headerbgcolor=blue!15,
    headertextcolor=blue!70!black,
    headerfont=\sffamily\bfseries
  }
\pxTable{Name, Age, City}
        {Alice, 20, Barcelona}
        {Bob, 22, Madrid}
\end{verbatim}

\pxTABsetup
  {
    style=header,
    headerbgcolor=blue!15,
    headertextcolor=blue!70!black,
    headerfont=\sffamily\bfseries
  }
\pxTable{Name, Age, City}
        {Alice, 20, Barcelona}
        {Bob, 22, Madrid}

\subsection{Striped}

\begin{verbatim}
\pxTABsetup
  {
    style=striped,
    altrowbgcolor=gray!10,
    altrowtextcolor=black
  }
\pxTable{Item, Quantity}
        {Apples, 3}
        {Oranges, 5}
        {Bananas, 2}

\end{verbatim}

\pxTABsetup
  {
    style=striped,
    altrowbgcolor=gray!10,
    altrowtextcolor=black
  }
\pxTable{Item, Quantity}
        {Apples, 3}
        {Oranges, 5}
        {Bananas, 2}

\subsection{Boxed}

\begin{verbatim}
\pxTABsetup{style=boxed}
\pxTableFromList
  {Code, Description;
   PX1, First code;
   PX2, Second code}
\end{verbatim}

\pxTABsetup{style=boxed}
\pxTableFromList
  {Code, Description;
   PX1, First code;
   PX2, Second code}

\subsection{Grid}

\begin{verbatim}
\pxTABsetup{style=grid}
\pxTableMatrix{2}{3}
  {m11,m12,m13,
   m21,m22,m23}
\end{verbatim}

\pxTABsetup{style=grid}
\pxTableMatrix{2}{3}
  {m11,m12,m13,
   m21,m22,m23}

\subsection{Composed header styles}

The composed styles combine the header row styling with the boxed,
grid, or striped variants so that common table layouts can be selected
with a single option.  The following examples use the same small
dataset.

\subsubsection*{\texttt{style=header-boxed}}

\begin{verbatim}
\pxTABsetup
  {
    style=header-boxed,
    headerbgcolor=blue!15,
    headertextcolor=blue!70!black,
    headerfont=\sffamily\bfseries,
    tabwidth=.8\linewidth
  }
\pxTable{Name, Age, City}
        {Alice, 20, Barcelona}
        {Bob, 22, Madrid}
\end{verbatim}

\pxTABsetup
  {
    style=header-boxed,
    headerbgcolor=blue!15,
    headertextcolor=blue!70!black,
    headerfont=\sffamily\bfseries,
    tabwidth=.8\linewidth
  }
\pxTable{Name, Age, City}
        {Alice, 20, Barcelona}
        {Bob, 22, Madrid}

\subsubsection*{\texttt{style=header-grid}}

\begin{verbatim}
\pxTABsetup
  {
    style=header-grid,
    headerbgcolor=blue!15,
    headertextcolor=blue!70!black,
    headerfont=\sffamily\bfseries,
    tabwidth=.8\linewidth
  }
\pxTable{Name, Age, City}
        {Alice, 20, Barcelona}
        {Bob, 22, Madrid}
\end{verbatim}

\pxTABsetup
  {
    style=header-grid,
    headerbgcolor=blue!15,
    headertextcolor=blue!70!black,
    headerfont=\sffamily\bfseries,
    tabwidth=.8\linewidth
  }
\pxTable{Name, Age, City}
        {Alice, 20, Barcelona}
        {Bob, 22, Madrid}

\subsubsection*{\texttt{style=header-striped}}

\begin{verbatim}
\pxTABsetup
  {
    style=header-striped,
    headerbgcolor=blue!15,
    headertextcolor=blue!70!black,
    headerfont=\sffamily\bfseries,
    altrowbgcolor=gray!10,
    altrowtextcolor=black,
    tabwidth=.8\linewidth
  }
\pxTable{Name, Age, City}
        {Alice, 20, Barcelona}
        {Bob, 22, Madrid}
        {Clara, 23, Girona}
\end{verbatim}

\pxTABsetup
  {
    style=header-striped,
    headerbgcolor=blue!15,
    headertextcolor=blue!70!black,
    headerfont=\sffamily\bfseries,
    altrowbgcolor=gray!10,
    altrowtextcolor=black,
    tabwidth=.8\linewidth
  }
\pxTable{Name, Age, City}
        {Alice, 20, Barcelona}
        {Bob, 22, Madrid}
        {Clara, 23, Girona}

\subsection{Centered and floating tables}

The new \texttt{centering} and float-related keys can be combined with
any of the predefined styles.  The following example uses
\texttt{style=header} with global centering:

\begin{verbatim}
\pxTABsetup
  {
    style=header,
    headerbgcolor=blue!15,
    headertextcolor=blue!70!black,
    headerfont=\sffamily\bfseries,
    centering=true
  }
\pxTable{Name, Age, City}
        {Alice, 20, Barcelona}
        {Bob, 22, Madrid}
\end{verbatim}

\pxTABsetup
  {
    style=header,
    headerbgcolor=blue!15,
    headertextcolor=blue!70!black,
    headerfont=\sffamily\bfseries,
    centering=true
  }
\pxTable{Name, Age, City}
        {Alice, 20, Barcelona}
        {Bob, 22, Madrid}

A float table with caption and label can be created by passing the new
keys in the options argument:

\begin{verbatim}
\pxTable[
  style=header,
  caption={Sample grades},
  label=tab:grades,
  float=htbp
]{Name, Grade}
 {Alice, 9}
 {Bob, 7}
\end{verbatim}

\pxTable[
  style=header,
  caption={Sample grades},
  label=tab:grades,
  float=htbp
]{Name, Grade}
 {Alice, 9}
 {Bob, 7}

As usual, the table can later be referenced with
\verb|Table~\ref{tab:grades}| in the surrounding text.

\section{Command reference}

This section summarises the user-facing commands provided by
\texttt{pxTAB}.  All commands respect the key settings described
above.

\subsection{\texttt{\textbackslash pxTABsetup}}

\begin{verbatim}
\pxTABsetup{<key>=<value>,...}
\end{verbatim}

Sets package-wide defaults in the \texttt{pxTAB} key family.
Usual \LaTeX\ scoping rules apply, so placing a setup call inside a
group affects only the tables in that group.

\subsection{\texttt{\textbackslash pxTable}}

\begin{verbatim}
\pxTable[<options>]{<row1>}{<row2>}...{<rowN>}
\end{verbatim}

Builds a table from braced row arguments.  Each \verb|<row>| is a
comma-separated list of cell contents.  A typical example is:

\begin{verbatim}
\pxTABsetup{style=header}
\pxTable{Name, Age, City}
        {Alice, 20, Barcelona}
        {Bob, 22, Madrid}
\end{verbatim}

\subsection{\texttt{\textbackslash pxTableFromList}}

\begin{verbatim}
\pxTableFromList[<options>]{<data>}
\end{verbatim}

Builds a table from a single list-style string.  By default rows are
separated by semicolons (\verb|;|) and cells by commas (\verb|,|):

\begin{verbatim}
\pxTableFromList
  {Name, Age, City;
   Alice, 20, Barcelona;
   Bob, 22, Madrid}
\end{verbatim}

The separators can be customised with the \texttt{rowsep} and
\texttt{cellsep} keys.

\subsection{\texttt{\textbackslash pxTableRow}}

\begin{verbatim}
\pxTableRow[<options>]{<row>}
\end{verbatim}

Convenience wrapper for single-row tables (headers or small data
rows).  It behaves as:

\begin{verbatim}
\pxTableRow[<opts>]{<row>} == \pxTable[<opts>]{<row>}
\end{verbatim}

\subsection{\texttt{\textbackslash pxTableHeadBody}}

\begin{verbatim}
\pxTableHeadBody[<options>]{<header>}{<body-list>}
\end{verbatim}

The header row is given explicitly as a comma-separated list, while
the body is provided as list-style data using the current separators:

\begin{verbatim}
\pxTableHeadBody
  {Name, Role}
  {Alice, Dev; Bob, QA; Carol, PO}
\end{verbatim}

\subsection{\texttt{\textbackslash pxTableHeadBodyRows}}

\begin{verbatim}
\pxTableHeadBodyRows[<options>]
  {<header>}{<row1>}{<row2>}...
\end{verbatim}

Variant of the previous command where both header and body rows are
given as separate braced arguments:

\begin{verbatim}
\pxTableHeadBodyRows
  {Name, Score}
  {Anna, 8}
  {Marc, 9}
  {Joan, 10}
\end{verbatim}

\subsection{\texttt{\textbackslash pxTableMatrix}}

\begin{verbatim}
\pxTableMatrix[<options>]{<rows>}{<cols>}{<flat-cells>}
\end{verbatim}

Builds a table from a flat list of cells, filled row by row.  Missing
cells are left blank:

\begin{verbatim}
\pxTableMatrix{2}{3}
  {m11,m12,m13,
   m21,m22,m23}
\end{verbatim}

\subsection{Quick reference}

The following table summarises the most common commands and their
arguments; all of them honour the global and per-table keys described in
the configuration section.

\begin{center}
\begin{tabular}{llp{.45\linewidth}}
\textbf{Command} & \textbf{Arguments} & \textbf{Typical use}\\\hline
\verb|\pxTABsetup| & \verb|{<key>=<value>,...}| &
  Set global defaults (style, width, centering, colours, float keys).\\
\verb|\pxTable| &
  \verb|[<options>]{<row1>}{...}{<rowN>}| &
  Small to medium tables entered as explicit rows.  Supports
  \texttt{centering}, \texttt{caption}, \texttt{label}, \texttt{float}.\\
\verb|\pxTableFromList| &
  \verb|[<options>]{<data>}| &
  Tables given as a single list string separated by \texttt{rowsep} and
  \texttt{cellsep}.\\
\verb|\pxTableRow| &
  \verb|[<options>]{<row>}| &
  Single-row tables (headers or quick inline tables).\\
\verb|\pxTableHeadBody| &
  \verb|[<options>]{<header>}{<body-list>}| &
  Head/body split with the body given as list-style data.\\
\verb|\pxTableHeadBodyRows| &
  \verb|[<options>]{<header>}{<row1>}{...}| &
  Head/body split where both header and body rows are given explicitly.\\
\verb|\pxTableMatrix| &
  \verb|[<options>]{<rows>}{<cols>}{<flat-cells>}| &
  Matrix-style data where a flat cell list is reshaped into a grid.\\
\end{tabular}
\end{center}

The float-related keys \texttt{caption}, \texttt{label}, and
\texttt{float}, together with \texttt{centering}, can be provided either
globally via \verb|\pxTABsetup| or per table in the \verb|[<options>]|
argument.

\section{Advanced usage}

\subsection{Best practices}

\begin{itemize}
  \item Use \texttt{style=plain} for compact, in-text tables where visual
    emphasis is not required; \texttt{style=header} works well for small
    datasets with a clear header row.
  \item For longer datasets, \texttt{style=striped} improves readability,
    while \texttt{boxed} and \texttt{grid} are suitable when a strong
    tabular structure is desired (summaries, rubrics, code lists).
  \item Prefer the natural-width default together with
    \texttt{centering=true} when the table is narrower than the text
    block; this keeps columns tight while visually centring the table.
  \item Use \verb|width=\linewidth| (or a fraction thereof) when the
    table should align with surrounding paragraphs; in this case
    \texttt{centering} is optional but harmless.
  \item When using floats, set both \texttt{caption} and \texttt{label}
    so that tables integrate with the document’s list of tables and can
    be referenced with \verb|\ref|.
\end{itemize}

\subsection{Explicit column specifications}

When \texttt{colspec} is set, its value is used directly as the
\texttt{tabular} column specification:

\begin{verbatim}
\pxTABsetup{style=plain,colspec=ll}
\pxTable{Col 1, Col 2}
        {Left, also left}
\end{verbatim}

Simple \texttt{l}/\texttt{c}/\texttt{r} specifications are checked
against the inferred column count; on mismatch a warning is issued
and an automatic specification is used.

\subsection{Width control}

The \texttt{width}/\texttt{tabwidth} keys allow fixed-width tables:

\begin{verbatim}
\pxTABsetup{style=plain,width=\linewidth,align=l}
\pxTable{A, B, C}{1, 2, 3}{4, 5, 6}
\end{verbatim}

Non-\texttt{natural} widths use a \texttt{tabular*} layout with
automatic extra column separation so that the table fits the requested
width.

\subsection{Font size and row spacing}

\begin{verbatim}
\pxTable[fontsize=\small,rowstretch=1.05,style=header]
        {Name, Notes}
        {Alice, Compact rows with smaller text}
        {Bob, This checks per-table sizing controls}
\end{verbatim}

The \texttt{fontsize}/\texttt{size} and \texttt{arraystretch}/%
\texttt{rowstretch} keys apply only to the current table.

\subsection{Troubleshooting}

\begin{itemize}
  \item \emph{Empty or missing tables.}  If the first row has no cells,
    \texttt{pxTAB} skips the table and issues a warning.  Check that
    your header row is non-empty and that the separators
    (\texttt{rowsep}, \texttt{cellsep}) are correct.
  \item \emph{Column specification mismatch.}  When \texttt{colspec} is
    given but its length does not match the inferred column count, an
    automatic specification is used instead and a warning is written to
    the log.  Adjust \texttt{colspec} or let \texttt{pxTAB} derive it
    from \texttt{align}.
  \item \emph{Float placement warnings.}  The \texttt{float} key is
    passed directly to \verb|\begin{table}[...]|.  If you use an
    unsupported placement specifier, \LaTeX\ may report an ``Unknown
    float option'' warning; choose a standard value such as
    \texttt{h}, \texttt{t}, or \texttt{htbp}.
  \item \emph{Missing caption or label.}  A float without \texttt{label}
    still typesets correctly but cannot be referenced with
    \verb|\ref|.  Conversely, a \texttt{label} without \texttt{caption}
    is allowed; the table is still wrapped in a \texttt{table}
    environment so that the label has a proper float context.
\end{itemize}

\section{Implementation notes}

This subsection describes, at a high level, how \texttt{pxTAB} processes
table data and renders it.  It is intended for maintainers and advanced
users who need to understand the internal pipeline.

\subsection*{Row parsing and trailing empty trimming}

Input data (from \verb|\pxTable|, \verb|\pxTableFromList|, or the
head/body helpers) is stored as a sequence of rows in
\verb|\l__pxtab_rows_seq|.  For list-style input, rows are split using
the \texttt{rowsep} separator (default \texttt{;}), and cells within each
row using \texttt{cellsep} (default \texttt{,}).

Before rendering, \verb|\__pxtab_trim_trailing_empty_rows:| recursively
removes trailing rows that have no non-blank cells.  A row is considered
empty if every cell, after splitting by \texttt{cellsep}, is blank.
This ensures that trailing semicolons (e.g.\ \texttt{A,B;C,D;}) do not
produce extra blank rows in the output.

\subsection*{Column count and colspec derivation}

The column count is taken from the \emph{first} row only: it is the
number of cells after splitting by \texttt{cellsep}.  If the first row
is empty, a warning is emitted and the table is skipped.

The column specification (\texttt{colspec}) is determined as follows:
\begin{itemize}
  \item If the user provided \texttt{colspec} and it is a simple
    \texttt{l}/\texttt{c}/\texttt{r} string whose length matches the
    inferred column count, it is used as-is.
  \item If the provided \texttt{colspec} length does not match, a
    \texttt{colspec-mismatch} warning is issued and an automatic
    colspec is built from the \texttt{align} key.
  \item If no \texttt{colspec} was given, one is built automatically
    by repeating the \texttt{align} character (default \texttt{c}) for
    each column.
\end{itemize}

For \texttt{boxed} and \texttt{grid} styles, vertical rules are added
around the colspec (e.g.\ \texttt{|c|c|c|}).

\subsection*{Style renderers and shared helpers}

Each style has a dedicated renderer
(\verb|\__pxtab_render_table_plain:n|, \verb|header:n|,
\verb|striped:n|, \verb|boxed:n|, \verb|grid:n|).  All renderers share:

\begin{itemize}
  \item \verb|\__pxtab_append_row_to_body:nnnn| --- appends a row of
    cells to the body token list, with optional text colour and font.
  \item \verb|\__pxtab_append_row_styled_to_body:nnnnn| --- same, plus
    optional \verb|\rowcolor| for background.
  \item \verb|\__pxtab_emit_tabular_with_body:N| --- applies visual
    setup (\texttt{tabcolsep}, \texttt{arraystretch}, \texttt{rulecolor},
    etc.) and emits the \texttt{tabular} or \texttt{tabular*} with the
    accumulated body.
\end{itemize}

Style-specific behaviour:
\begin{itemize}
  \item \texttt{plain}: all rows use body font; no row colours or rules.
  \item \texttt{header}: first row gets \texttt{headerbgcolor},
    \texttt{headertextcolor}, \texttt{headerfont}; body rows use body font.
  \item \texttt{striped}: odd rows plain; even rows get
    \texttt{altrowbgcolor} and \texttt{altrowtextcolor}.
  \item \texttt{boxed}: colspec wrapped with \texttt{|}; \texttt{\textbackslash hline}
    before first row and after last row.
  \item \texttt{grid}: colspec wrapped with \texttt{|}; \texttt{\textbackslash hline}
    before and after every row.
  \item \texttt{header-boxed}: header-style first row (colours and font)
    with boxed frame rules as in \texttt{boxed}.
  \item \texttt{header-grid}: header-style first row together with full
    grid rules as in \texttt{grid}.
  \item \texttt{header-striped}: header-style first row; subsequent body
    rows alternate plain and striped colours like \texttt{striped}, with
    the striping parity starting at the first body row.
\end{itemize}

\section{Integration with P3CTeX}

When using the \texttt{P3CTeX} document class and its core package
\texttt{pxCORE}, the \texttt{pxTAB} package is loaded as part of the
default configuration so that table commands are available without
extra setup.  The example document \texttt{P3CTeX-example.tex} in the
repository includes several small \texttt{pxTAB} tables alongside the
other P3CTeX modules.

\end{document}

